%% kapa: exam=1 round=28 year=2017 period=1 form=A subject=real_estate subject_ko=부동산학원론 seq=22 number=102
%% source: Q-Net 감정평가사 기출문제 (cst00309) artlSeq=5199603 file=extracted/1008/5199603/2017년 제28회 감정평가사 1차시험 문제지/1교시 A형.pdf
%% conversion: kordoc 4.24.51 -> tools/gen1.py (auto); review=proofread-sonnet; flags=none
\begin{question}{102}
부동산평가활동에서 부동산가격의 원칙에 관한 설명으로 옳지 않은 것은?
\begin{choices}
  \item 기여의 원칙이란 부동산가격이 대상부동산의 각 구성요소가 기여하는 정도의 합으로 결정된다는 것을 말한다.
  \item 최유효이용의 원칙이란 객관적으로 보아 양식과 통상의 이용능력을 지닌 사람이 대상토지를 합법적이고 합리적이며 최고최선의 방법으로 이용하는 것을 말한다.
  \item 변동의 원칙이란 가치형성요인이 시간의 흐름에 따라 지속적으로 변화함으로써 부동산가격도 변화한다는 것을 말한다.
  \item 적합의 원칙이란 부동산의 유용성이 최고도로 발휘되기 위해서는 부동산구성요소의 결합에 균형이 있어야 한다는 것을 말한다.
  \item 예측의 원칙이란 평가활동에서 가치형성요인의 변동추이 또는 동향을 주시해야 한다는 것을 말한다.
\end{choices}
\end{question}
\answer{4}
